% ==============================================================================
% 2. PROFIL KELUARGA DAN MATA PENCAHARIAN
% ==============================================================================

\section{Profil Keluarga dan Mata Pencaharian}

Sektor mata pencaharian utama keluarga di Desa Popontolen terbagi atas sektor pertanian dan sektor non-pertanian. Secara umum, mayoritas keluarga di desa ini mengandalkan sektor non-pertanian (seperti perdagangan kelontong, kedai pangan, jasa, wiraswasta, maupun pegawai) dengan total \textbf{272 keluarga} pada tahun 2025.

Sektor pertanian dan perkebunan (khususnya kelapa dan tanaman pangan) merupakan pilar penting di wilayah tertentu, dengan total \textbf{179 keluarga petani}. Wilayah dengan konsentrasi keluarga petani terbesar berada di \textbf{Jaga 7} (67 KK petani), disusul oleh \textbf{Jaga 6} (30 KK petani), \textbf{Jaga 2} (27 KK petani), dan \textbf{Jaga 4} (24 KK petani).

\vspace{0.3cm}

\begin{table}[htbp]
\centering
\caption{Jumlah Keluarga Berdasarkan Sektor Pekerjaan Utama (2024--2025)}
\label{tab:pekerjaan_keluarga}
\rowcolors{3}{bpslight}{white}
\begin{tabularx}{\textwidth}{l rrr rrr}
\toprule
\rowcolor{bpsnavy}
\textcolor{white}{\textbf{Wilayah Jaga}} & 
\textcolor{white}{\textbf{Total KK (24)}} & 
\textcolor{white}{\textbf{Tani (24)}} & 
\textcolor{white}{\textbf{Non-T (24)}} & 
\textcolor{white}{\textbf{Total KK (25)}} & 
\textcolor{white}{\textbf{Tani (25)}} & 
\textcolor{white}{\textbf{Non-T (25)}} \\
\midrule
Jaga 1 & 52 & 8 & 44 & 52 & 8 & 44 \\
Jaga 2 & 62 & 27 & 35 & 62 & 27 & 35 \\
Jaga 3 & 50 & 12 & 28 & 50 & 12 & 28 \\
Jaga 4 & 75 & 24 & 51 & 79 & 24 & 55 \\
Jaga 5 & 56 & 11 & 0 & 56 & 11 & 0 \\
Jaga 6 & 78 & 30 & 18 & 74 & 30 & 16 \\
Jaga 7 & 161 & 67 & 94 & 160 & 67 & 94 \\
\midrule
\rowcolor{bpsborder!40}
\textbf{TOTAL DESA} & \textbf{534} & \textbf{179} & \textbf{270} & \textbf{533} & \textbf{179} & \textbf{272} \\
\bottomrule
\end{tabularx}
\end{table}

Berdasarkan proporsi pekerjaan di atas, Pemerintah Desa Popontolen menitikberatkan kebijakan pembangunan pada penguatan rantai nilai perkebunan kelapa dan pertanian pangan terpadu, sekaligus pembinaan manajemen dan akses permodalan bagi usaha mikro perdagangan di setiap wilayah Jaga.
